problem:  
Let \((M^n,g)\) be a compact Riemannian manifold and denote by \(\lambda_{1,p}(M,g)\) the first nonzero eigenvalue of the \(p\)-Laplacian  
\[
\Delta_p u = \operatorname{div}(|\nabla u|^{p-2}\nabla u)
\]
with homogeneous Neumann boundary conditions (or on the closed manifold).  
It is known that  
\[
\lambda_{1,p}(M,g)^{1/p} \to \frac{2}{\operatorname{diam}(M,g)} \quad \text{as } p\to\infty.
\]

Define the *renormalized deviation*  
\[
\delta_p(M,g) := p\left[\left(\frac{\operatorname{diam}(M,g)}{2}\right)\lambda_{1,p}(M,g)^{1/p} - 1\right].
\]

**Open problem:**  
Does \(\delta_p(M,g)\) have a finite limit as \(p \to \infty\), and if so, is this limit determined by curvature invariants or geometric features of \((M,g)\)?  
More specifically:  
1. Is the sign or finiteness of \(\lim_{p\to\infty} \delta_p(M,g)\) influenced by lower Ricci curvature bounds, e.g., \(\operatorname{Ric}_g \ge -K\)?  
2. Can one give an explicit geometric interpretation of the first correction term in the asymptotic expansion
\[
\lambda_{1,p}(M,g)^{1/p} = \frac{2}{\operatorname{diam}(M,g)}\Big(1 + \frac{A(M,g)}{p} + o\Big(\frac{1}{p}\Big)\Big),
\]
for some geometric functional \(A(M,g)\) possibly related to curvature, boundary structure, or the distribution of midpoints of geodesic pairs realizing the diameter?  
3. What is the relationship between the asymptotic profile of \(p\)-eigenfunctions and such higher-order corrections—does curvature manifest in their limiting concentration patterns as \(p\to\infty\)?  

Why is it a "good" problem:  
This problem addresses the *second-order asymptotics* of nonlinear eigenvalues, a frontier region beyond the classical \(p\to\infty\) limit known to yield only diameter information. It effectively asks whether finer geometric information—such as curvature, topology, or cut-locus structure—can be extracted from the rate at which \(\lambda_{1,p}^{1/p}\) approaches its limit.  

The problem is "good" for several reasons:  

1. **Conceptual depth:** It moves beyond the well-understood diameter dependence of \(\Lambda_\infty\), probing whether nonlinear spectral data encode additional geometric subtleties invisible in the limit itself.  

2. **Cross-field bridge:** It naturally connects PDE asymptotics, Riemannian geometry, and geometric measure theory—analyzing how curvature affects viscosity-limit eigenfunctions of the infinity-Laplacian.  

3. **Simplicity with depth:** The statement involves only one familiar operator (\(p\)-Laplacian), one metric invariant (diameter), and a simple asymptotic question, yet its resolution may require developing new geometric perturbation theories for nonlinearly elliptic spectra.  

A success in this direction would initiate a theory of *nonlinear spectral asymptotics* paralleling semiclassical analysis, revealing how manifold geometry is subtly encoded in the nonlinear spectrum beyond the known diameter constraint.